\documentclass[12pt,a4paper]{article}

% ---- Encoding & Font ----
\usepackage{iftex}
\ifPDFTeX
  \usepackage[utf8]{inputenc}
  \usepackage[T1]{fontenc}
\else
  \usepackage{fontspec}  % XeTeX/LuaTeX (tectonic): Unicode nativo (·, —, ¿, ª, §)
\fi
\usepackage[spanish]{babel}
\usepackage{csquotes}

% ---- Page layout ----
\usepackage[top=2.5cm, bottom=2.5cm, left=2.5cm, right=2.5cm]{geometry}
\usepackage{setspace}
\onehalfspacing

% ---- Math ----
\usepackage{amsmath, amssymb, amsthm}
\usepackage{mathtools}
\usepackage{physics}

% ---- Graphics ----
\usepackage{graphicx}
\usepackage{tikz}
\usepackage{float}
\usepackage{caption}

% ---- Tables ----
\usepackage{array}
\usepackage{booktabs}
\usepackage{tabularx}
\usepackage{colortbl}

% ---- Colours ----
\usepackage{xcolor}
\definecolor{notebg}{HTML}{FFF3CD}
\definecolor{noteborder}{HTML}{FFC107}
\definecolor{boxred}{HTML}{F8D7DA}
\definecolor{boxredborder}{HTML}{F5C6CB}

% ---- Boxes ----
\usepackage{mdframed}
\usepackage{tcolorbox}
\tcbuselibrary{most}

\newtcolorbox{keybox}[1][]{
  colback=blue!5,
  colframe=blue!70!black,
  arc=4pt,
  boxrule=1pt,
  fonttitle=\bfseries,
  title={#1}
}

\newtcolorbox{warnbox}[1][]{
  colback=boxred,
  colframe=boxredborder,
  coltitle=black!75,
  arc=4pt,
  boxrule=1pt,
  fonttitle=\bfseries,
  title={#1}
}

\newtcolorbox{notebox}[1][]{
  colback=notebg,
  colframe=noteborder,
  coltitle=black!75,
  arc=4pt,
  boxrule=1pt,
  fonttitle=\bfseries,
  title={#1}
}

% ---- Diagramas (gráficas) ----
\usepackage{pgfplots}
\pgfplotsset{compat=1.18}
\usetikzlibrary{babel}  % babel español activa `>`; esto evita romper las flechas `->`
% courses/CDIV2017/Clases/assets/tikzstyles.tex
% ---------------------------------------------------------------------------
% Estilos compartidos para los diagramas del curso Cálculo 1 (CDIV2017).
% Fuente única del "look" visual (colores, grosores, escalas) para que todas
% las figuras (TikZ / pgfplots: conjuntos, rectas numéricas, gráficas de
% funciones) sean consistentes entre clases.
%
% Uso: la clase debe cargar antes `tikz` y `pgfplots` (bloque §6.3 del
%      CLAUDE.md del curso), y luego
%      % courses/CDIV2017/Clases/assets/tikzstyles.tex
% ---------------------------------------------------------------------------
% Estilos compartidos para los diagramas del curso Cálculo 1 (CDIV2017).
% Fuente única del "look" visual (colores, grosores, escalas) para que todas
% las figuras (TikZ / pgfplots: conjuntos, rectas numéricas, gráficas de
% funciones) sean consistentes entre clases.
%
% Uso: la clase debe cargar antes `tikz` y `pgfplots` (bloque §6.3 del
%      CLAUDE.md del curso), y luego
%      % courses/CDIV2017/Clases/assets/tikzstyles.tex
% ---------------------------------------------------------------------------
% Estilos compartidos para los diagramas del curso Cálculo 1 (CDIV2017).
% Fuente única del "look" visual (colores, grosores, escalas) para que todas
% las figuras (TikZ / pgfplots: conjuntos, rectas numéricas, gráficas de
% funciones) sean consistentes entre clases.
%
% Uso: la clase debe cargar antes `tikz` y `pgfplots` (bloque §6.3 del
%      CLAUDE.md del curso), y luego
%      \input{../assets/tikzstyles.tex}
% (compilar desde el directorio de la clase para que la ruta relativa resuelva).
%
% Paleta propia de este curso (no se reusa la de courses/Fisica3): mismos
% tonos base (azul/rojo/ámbar) para alinear con keybox/notebox/warnbox, pero
% archivo independiente, sin dependencia cruzada entre cursos.
% ---------------------------------------------------------------------------

% ---- Paleta (alineada con las cajas del preámbulo: keybox/notebox/warnbox) ----
\definecolor{figblue}{HTML}{1F4E79}   % azul principal (líneas, keybox, conjuntos)
\definecolor{figred}{HTML}{B03A2E}    % rojo de énfasis (warnbox, región resaltada)
\definecolor{figamber}{HTML}{B8860B}  % ámbar (notebox)
\definecolor{figgray}{HTML}{555555}   % auxiliares, ejes, guías, universo

% ---- TikZ: librerías y estilos reutilizables ----
% Nota: las puntas se declaran como `-{Latex[...]}` y no como `->`. La forma
% con llaves NO contiene un `>` literal, así que es inmune al `>` activo que
% introduce babel español (de todos modos se carga \usetikzlibrary{babel} en
% el bloque §6.3 para cubrir cualquier `->` suelto).
\usetikzlibrary{arrows.meta,calc,angles,quotes,patterns}

\tikzset{
  figline/.style={figblue, line width=0.9pt},
  figaux/.style={figgray, dashed, line width=0.6pt},
  figaxis/.style={figgray, line width=0.7pt, -{Latex[length=2.2mm,width=1.5mm]}},
  % conjuntos (diagramas de Venn / patatoidales)
  figset/.style={figline, fill=figblue!6},
  figsetB/.style={figred, line width=0.9pt, fill=figred!6},
  figsetC/.style={figamber, line width=0.9pt, fill=figamber!10},
  figuniverso/.style={figgray, line width=0.7pt},
  % puntos de la recta / conjuntos numéricos
  figpto/.style={circle, inner sep=0pt, minimum size=5pt, draw=none},
  figptoabierto/.style={circle, inner sep=0pt, minimum size=5pt, draw=figblue, line width=0.9pt, fill=white},
  figptocerrado/.style={circle, inner sep=0pt, minimum size=5pt, draw=figblue, line width=0.9pt, fill=figblue},
  figptoY/.style={figpto, fill=figblue},
  figptoZ/.style={figpto, fill=figred},
  % segmentos gruesos para intervalos sobre la recta
  figintervalo/.style={figblue, line width=2.2pt},
  % vectores/flechas de transformación (traslación, dilatación)
  figvec/.style={figred, line width=1pt, -{Latex[length=2.2mm,width=1.6mm]}},
  % etiquetas
  figlbl/.style={font=\small},
  figlblaux/.style={font=\footnotesize, figgray},
}

% ---- pgfplots: estilo base de las gráficas de funciones ----
\pgfplotsset{
  calcfig/.style={
    width=10cm, height=6cm,
    axis lines=middle,
    xlabel style={font=\small, at={(axis description cs:1,0.5)}, anchor=west},
    ylabel style={font=\small, at={(axis description cs:0.5,1)}, anchor=south},
    tick label style={font=\footnotesize},
    every axis plot/.append style={line width=1.1pt, figblue},
    grid=none,
    legend cell align=left,
    legend style={font=\footnotesize, draw=figgray!40},
    clip=false,
  },
}

% (compilar desde el directorio de la clase para que la ruta relativa resuelva).
%
% Paleta propia de este curso (no se reusa la de courses/Fisica3): mismos
% tonos base (azul/rojo/ámbar) para alinear con keybox/notebox/warnbox, pero
% archivo independiente, sin dependencia cruzada entre cursos.
% ---------------------------------------------------------------------------

% ---- Paleta (alineada con las cajas del preámbulo: keybox/notebox/warnbox) ----
\definecolor{figblue}{HTML}{1F4E79}   % azul principal (líneas, keybox, conjuntos)
\definecolor{figred}{HTML}{B03A2E}    % rojo de énfasis (warnbox, región resaltada)
\definecolor{figamber}{HTML}{B8860B}  % ámbar (notebox)
\definecolor{figgray}{HTML}{555555}   % auxiliares, ejes, guías, universo

% ---- TikZ: librerías y estilos reutilizables ----
% Nota: las puntas se declaran como `-{Latex[...]}` y no como `->`. La forma
% con llaves NO contiene un `>` literal, así que es inmune al `>` activo que
% introduce babel español (de todos modos se carga \usetikzlibrary{babel} en
% el bloque §6.3 para cubrir cualquier `->` suelto).
\usetikzlibrary{arrows.meta,calc,angles,quotes,patterns}

\tikzset{
  figline/.style={figblue, line width=0.9pt},
  figaux/.style={figgray, dashed, line width=0.6pt},
  figaxis/.style={figgray, line width=0.7pt, -{Latex[length=2.2mm,width=1.5mm]}},
  % conjuntos (diagramas de Venn / patatoidales)
  figset/.style={figline, fill=figblue!6},
  figsetB/.style={figred, line width=0.9pt, fill=figred!6},
  figsetC/.style={figamber, line width=0.9pt, fill=figamber!10},
  figuniverso/.style={figgray, line width=0.7pt},
  % puntos de la recta / conjuntos numéricos
  figpto/.style={circle, inner sep=0pt, minimum size=5pt, draw=none},
  figptoabierto/.style={circle, inner sep=0pt, minimum size=5pt, draw=figblue, line width=0.9pt, fill=white},
  figptocerrado/.style={circle, inner sep=0pt, minimum size=5pt, draw=figblue, line width=0.9pt, fill=figblue},
  figptoY/.style={figpto, fill=figblue},
  figptoZ/.style={figpto, fill=figred},
  % segmentos gruesos para intervalos sobre la recta
  figintervalo/.style={figblue, line width=2.2pt},
  % vectores/flechas de transformación (traslación, dilatación)
  figvec/.style={figred, line width=1pt, -{Latex[length=2.2mm,width=1.6mm]}},
  % etiquetas
  figlbl/.style={font=\small},
  figlblaux/.style={font=\footnotesize, figgray},
}

% ---- pgfplots: estilo base de las gráficas de funciones ----
\pgfplotsset{
  calcfig/.style={
    width=10cm, height=6cm,
    axis lines=middle,
    xlabel style={font=\small, at={(axis description cs:1,0.5)}, anchor=west},
    ylabel style={font=\small, at={(axis description cs:0.5,1)}, anchor=south},
    tick label style={font=\footnotesize},
    every axis plot/.append style={line width=1.1pt, figblue},
    grid=none,
    legend cell align=left,
    legend style={font=\footnotesize, draw=figgray!40},
    clip=false,
  },
}

% (compilar desde el directorio de la clase para que la ruta relativa resuelva).
%
% Paleta propia de este curso (no se reusa la de courses/Fisica3): mismos
% tonos base (azul/rojo/ámbar) para alinear con keybox/notebox/warnbox, pero
% archivo independiente, sin dependencia cruzada entre cursos.
% ---------------------------------------------------------------------------

% ---- Paleta (alineada con las cajas del preámbulo: keybox/notebox/warnbox) ----
\definecolor{figblue}{HTML}{1F4E79}   % azul principal (líneas, keybox, conjuntos)
\definecolor{figred}{HTML}{B03A2E}    % rojo de énfasis (warnbox, región resaltada)
\definecolor{figamber}{HTML}{B8860B}  % ámbar (notebox)
\definecolor{figgray}{HTML}{555555}   % auxiliares, ejes, guías, universo

% ---- TikZ: librerías y estilos reutilizables ----
% Nota: las puntas se declaran como `-{Latex[...]}` y no como `->`. La forma
% con llaves NO contiene un `>` literal, así que es inmune al `>` activo que
% introduce babel español (de todos modos se carga \usetikzlibrary{babel} en
% el bloque §6.3 para cubrir cualquier `->` suelto).
\usetikzlibrary{arrows.meta,calc,angles,quotes,patterns}

\tikzset{
  figline/.style={figblue, line width=0.9pt},
  figaux/.style={figgray, dashed, line width=0.6pt},
  figaxis/.style={figgray, line width=0.7pt, -{Latex[length=2.2mm,width=1.5mm]}},
  % conjuntos (diagramas de Venn / patatoidales)
  figset/.style={figline, fill=figblue!6},
  figsetB/.style={figred, line width=0.9pt, fill=figred!6},
  figsetC/.style={figamber, line width=0.9pt, fill=figamber!10},
  figuniverso/.style={figgray, line width=0.7pt},
  % puntos de la recta / conjuntos numéricos
  figpto/.style={circle, inner sep=0pt, minimum size=5pt, draw=none},
  figptoabierto/.style={circle, inner sep=0pt, minimum size=5pt, draw=figblue, line width=0.9pt, fill=white},
  figptocerrado/.style={circle, inner sep=0pt, minimum size=5pt, draw=figblue, line width=0.9pt, fill=figblue},
  figptoY/.style={figpto, fill=figblue},
  figptoZ/.style={figpto, fill=figred},
  % segmentos gruesos para intervalos sobre la recta
  figintervalo/.style={figblue, line width=2.2pt},
  % vectores/flechas de transformación (traslación, dilatación)
  figvec/.style={figred, line width=1pt, -{Latex[length=2.2mm,width=1.6mm]}},
  % etiquetas
  figlbl/.style={font=\small},
  figlblaux/.style={font=\footnotesize, figgray},
}

% ---- pgfplots: estilo base de las gráficas de funciones ----
\pgfplotsset{
  calcfig/.style={
    width=10cm, height=6cm,
    axis lines=middle,
    xlabel style={font=\small, at={(axis description cs:1,0.5)}, anchor=west},
    ylabel style={font=\small, at={(axis description cs:0.5,1)}, anchor=south},
    tick label style={font=\footnotesize},
    every axis plot/.append style={line width=1.1pt, figblue},
    grid=none,
    legend cell align=left,
    legend style={font=\footnotesize, draw=figgray!40},
    clip=false,
  },
}


% ---- Hyperlinks & TOC ----
\usepackage[hidelinks]{hyperref}
\usepackage{bookmark}

% ---- Enumerate ----
\usepackage{enumitem}

% ---- Miscellaneous ----
\usepackage{icomma}
\usepackage{siunitx}
\usepackage{textcomp}
\usepackage[bottom]{footmisc}

% ---- Title ----
\title{\textbf{Clase 19: Límites --- ejemplos y contraejemplos} \\[0.3em]
        \large{La identidad, la raíz cuadrada, la negación de la definición y la función signo}}
\author{Transcripto y expandido --- Curso Cálculo 1 (OpenFING)}
\date{}

\begin{document}

\maketitle
\thispagestyle{empty}
\newpage

\tableofcontents
\newpage

\section{La definición de límite, repasada}

La clase abre recordando la definición dada la clase anterior, que el docente
califica sin rodeos como \textbf{la definición más difícil de todo el curso} y
la única que hay que memorizar si sólo se memoriza una cosa.

El marco es el de siempre en este capítulo: una función $f : I \to \mathbb{R}$
definida en un \textbf{intervalo $I$ de interior no vacío}, y un punto
$x_0 \in \bar I$, la \textbf{clausura} de $I$. Que el punto se tome en la
clausura y no en $I$ significa que $x_0$ puede estar adentro del intervalo o
ser uno de sus dos extremos, incluso si esos extremos no pertenecen a $I$. Es
un regalo gratis: la clausura agrega a lo sumo dos puntos más donde la noción
de límite tiene sentido, y si los extremos ya estaban adentro no cambia nada.

Dado además un número $L \in \mathbb{R}$, las dos notaciones
\[
  \lim_{x \to x_0} f(x) = L
  \qquad \text{y} \qquad
  f(x) \to L \ \text{ cuando } \ x \to x_0
\]
son \textbf{sinónimas}, y significan por definición:

\begin{keybox}[Definición: límite de una función en un punto]
\[
  \boxed{\ \forall \varepsilon > 0,\ \exists \delta_\varepsilon > 0,\
  \forall x \in I,\quad
  0 < |x - x_0| < \delta_\varepsilon \ \Longrightarrow\ |f(x) - L| < \varepsilon\ }
\]
\end{keybox}

Cada pedazo dice algo preciso. La condición $0 < |x - x_0|$ \textbf{no dice
nada más que $x \neq x_0$}: no hay ninguna sutileza escondida ahí. Junto con
$|x-x_0| < \delta$, lo que expresa es que $x$ pertenece al \textbf{entorno
reducido} de centro $x_0$ y radio $\delta$, y es reducido precisamente porque
el punto $x_0$ está prohibido. La conclusión, en cambio, usa un entorno
normal: $f(x)$ pertenece al entorno de centro $L$ y radio $\varepsilon$. En
notación de entornos:
\[
  x \in E^{*}_{\delta}(x_0) \ \Longrightarrow\ f(x) \in E_{\varepsilon}(L).
\]

Todo esto formaliza matemáticamente una idea intuitiva que conviene tener
siempre presente: \emph{$f(x)$ es arbitrariamente cercano a $L$ con tal de que
$x$ sea suficientemente cercano a $x_0$}.

\subsection{Un punto de lógica: los cuantificadores son símbolos}

El docente hace un paréntesis a raíz de observaciones que recibió sobre su
notación. Él escribe $\forall \varepsilon > 0,$ y $\exists \delta > 0,$ usando
una coma, que se lee ``tenemos que'' después de un para todo y ``tal que''
después de un existe. Muchos profesores ---en Francia y en Uruguay por igual---
escriben en cambio $\exists \delta > 0 \ /\ \dots$ o
$\exists \delta > 0 \ \text{tal que} \ \dots$, agregando el ``tal que'' al lado
del símbolo.

\begin{notebox}[Símbolo no es abreviatura]
Eso \textbf{está mal}, y la razón es una confusión entre \textbf{símbolo} y
\textbf{abreviatura}. El signo $\exists$ no es una estenografía de la palabra
``existe'': es un símbolo matemático, y contiene más palabras adentro. La
analogía es el signo de integral: $\int$ tampoco abrevia la palabra
``integral''. Si lo fuera, habría que escribir ``la $\int$ de $a$ a $b$ de
$f(x)\,dx$'', y no es así como se escribe.
\end{notebox}

Cuando se escribe en lenguaje natural, con palabras, el ``tal que'' \textbf{sí}
se restaura: ``sea $x$ tal que\dots'' es correcto, exactamente igual que decir
``la integral de $a$ a $b$ de $f$''. Lo que está mal es mezclar el símbolo con
la locución. Sobre el signo de puntuación: hay dos convenciones en lógica, o
bien paréntesis sin puntuación, o bien coma (o punto) sin paréntesis; en el
curso se usa la segunda, que es más cómoda.

\subsection{El juego entre el implementor y el usuario}

La dificultad real de la definición es la \textbf{alternancia de
cuantificadores}: un para todo, después un existe, después otro para todo. La
forma de entenderla es leerla como un juego entre dos jugadores, o ---para el
público informático de la sala--- como una \textbf{especificación de
software}, que tiene dos sentidos de lectura según de qué lado se esté.

\begin{center}
\begin{tabularx}{\textwidth}{@{}lXXX@{}}
\toprule
 & Elige $\varepsilon$ & Construye $\delta_\varepsilon$ & Elige el $x$ de prueba \\
\midrule
\textbf{Implementor} (demuestra que hay límite) & no, lo recibe & \textbf{sí} & no, lo recibe \\
\textbf{Usuario} (usa el límite como hipótesis) & \textbf{sí} & no, lo recibe & \textbf{sí} \\
\bottomrule
\end{tabularx}
\end{center}

El \textbf{implementor} es quien tiene que demostrar que la función tiene
límite. Desde su lugar, alguien de afuera le da un $\varepsilon$ sobre el cual
no tiene ningún control, y su trabajo es \textbf{construir} un $\delta$ en
función de ese $\varepsilon$. Por eso está bien ---y el docente lo hace
sistemáticamente--- escribir $\delta_\varepsilon$ con subíndice: \textbf{el
radio depende de la precisión}. Construir esa dependencia es literalmente
definir una función $\varepsilon \mapsto \delta_\varepsilon$, y en cada ejemplo
de la clase esa función se exhibe explícitamente. Después el usuario le da un
$x$ cualquiera para testear si el $\delta$ devuelto aguanta; por eso el
$\delta$ tiene que ser \textbf{genérico}, tiene que servir para todos los $x$ a
la vez.

El \textbf{usuario} tiene el punto de vista dual. Cuando la propiedad del
límite aparece como hipótesis de un teorema, ahora es él quien elige el
$\varepsilon$ que le conviene, recibe un $\delta_\varepsilon$ sobre el cual no
decide nada, y después elige el $x$ con el que va a explotar la información.

\begin{notebox}[La misma asimetría que en una función de un programa]
El implementor acepta el argumento sin control sobre su valor y tiene que
hacer algo con él; el usuario elige el argumento libremente pero no controla
lo que la función devuelve. Los conceptos de la informática ayudan mucho a
leer razonamientos matemáticos de este tipo.
\end{notebox}

\subsection{Lectura gráfica: banda de precisión y radio de seguridad}

El límite $L$ vive en el eje vertical (es un valor con el que se compara
$f(x)$), y $x_0$ en el horizontal. Un $\varepsilon$ dado define una
\textbf{banda horizontal de precisión} de semiancho $\varepsilon$ alrededor de
$L$. El juego consiste en hallar un $\delta$ \textbf{pequeño pero no nulo} tal
que la imagen de la banda vertical de radio $\delta$ alrededor de $x_0$ quede
íntegramente dentro de la banda horizontal. Un $\delta$ demasiado grande falla
porque parte de la gráfica se escapa por arriba o por abajo de la banda de
precisión; achicarlo lo arregla.

% Clase 19 --- El juego epsilon/delta y el entorno reducido (§1.3).
% El docente: "si un usuario me da un epsilon esto me define una banda
% horizontal, una banda de precisión... el juego es hallar un delta bastante
% pequeño, pero no nulo, tal que la imagen de la banda de radio delta esté
% adentro de la banda horizontal... yo no podría tomar delta acá, porque esta
% parte de la función se va afuera de la banda". Y además: "nunca se examina
% el valor de f en el punto x sub cero".
\begin{center}
\begin{tikzpicture}
\begin{axis}[
    calcfig,
    axis lines=left,
    width=12cm, height=7cm,
    domain=1.2:6.8, samples=100,
    xmin=0.8, xmax=7.6, ymin=1.4, ymax=5.4,
    xtick={4}, xticklabels={$x_0$},
    ytick={3}, yticklabels={$L$},
    xlabel={$x$}, ylabel={$y$},
    xlabel style={font=\small, at={(axis description cs:1,0)}, anchor=north},
    ylabel style={font=\small, at={(axis description cs:0,1)}, anchor=east},
  ]
  % banda horizontal de precisión: L +- epsilon
  \fill[figamber!18] (axis cs:0.8,2.4) rectangle (axis cs:7.05,3.6);
  % banda vertical que SÍ funciona: radio delta
  \fill[figblue!14] (axis cs:3.1,1.4) rectangle (axis cs:4.9,5.05);
  % radio demasiado grande: la imagen se escapa de la banda de precisión
  \draw[figred, dashed, line width=0.8pt] (axis cs:1.8,1.4) -- (axis cs:1.8,5.05);
  \draw[figred, dashed, line width=0.8pt] (axis cs:6.2,1.4) -- (axis cs:6.2,5.05);

  \addplot[figblue, line width=1.1pt] {3 + 0.6*(x-4) + 0.08*(x-4)^2};

  % guías del punto y del valor irrelevante f(x_0)
  \draw[figaux] (axis cs:0.8,3) -- (axis cs:4,3);
  \node[figptoabierto] at (axis cs:4,3) {};
  \node[figpto, fill=figred, minimum size=5.5pt] at (axis cs:4,4.6) {};
  \node[figlbl, figred, right=3pt] at (axis cs:4,4.6) {$f(x_0)$};

  % rótulos
  \node[figlbl, figamber, right] at (axis cs:7.08,3.42) {$L+\varepsilon$};
  \node[figlbl, figamber, right] at (axis cs:7.08,2.58) {$L-\varepsilon$};
  \node[figlblaux, align=center] at (axis cs:4,1.75)
    {radio $\delta_\varepsilon$\\ (sirve)};
  \node[figlbl, figred, align=center, anchor=south] at (axis cs:6.2,5.10)
    {radio demasiado grande};
\end{axis}
\end{tikzpicture}\\[0.5em]
{\small\itshape Una precisión $\varepsilon$ define la banda horizontal
$L\pm\varepsilon$; el juego consiste en devolver un radio $\delta_\varepsilon$
tal que la imagen de la banda vertical caiga entera adentro de la horizontal.
El radio punteado en rojo fracasa: parte de la gráfica se escapa. El valor
$f(x_0)$ puede estar en cualquier lado --- el entorno es \emph{reducido} y
nunca se lo mira.}
\end{center}


\begin{warnbox}[El límite nunca mira $f(x_0)$]
En este dibujo nunca se examina el valor $f(x_0)$. Ésa es exactamente la razón
de tomar entornos \textbf{reducidos}: la noción de límite intenta predecir lo
que pasa en $x_0$ \emph{sin mirar} lo que pasa en $x_0$. El valor $f(x_0)$
puede estar en cualquier otro lado del plano y el límite no cambia.
\end{warnbox}

\section{Continuidad}

Sobre la definición de límite se monta un caso particular, ya visto la clase
anterior. Ahora se exige $x_0 \in I$ ---\textbf{no} en la clausura---, porque
hay que poder hablar del valor $f(x_0)$:

\begin{keybox}[Definición: continuidad en un punto]
\[
  \boxed{\ f \text{ es continua en } x_0 \in I
  \iff \lim_{x \to x_0} f(x) \ \text{existe y vale} \ f(x_0)\ }
\]
\end{keybox}

Escribir $\lim_{x\to x_0} f(x) = f(x_0)$ afirma implícitamente \textbf{dos}
cosas: que el límite existe, y que coincide con la imagen.

\subsection{Los dos tipos de discontinuidad}

De ahí salen exactamente dos maneras de que una función \textbf{no} sea
continua en un punto:

\begin{enumerate}[label=\arabic*.]
  \item \textbf{Por ausencia de límite.} La función no tiene límite en $x_0$
        (oscila demasiado rápido, o los límites laterales difieren).
  \item \textbf{Porque el límite existe pero es distinto del valor.} Hay límite
        $L$, pero $f(x_0) \neq L$.
\end{enumerate}

El segundo caso es el más engañoso, porque nada falla ``cerca'' del punto: la
función se acerca ordenadamente a $L$ y el único problema es el valor aislado
que toma exactamente en $x_0$.

\subsection{Continuidad en un intervalo: cuatro cuantificadores}

\[
  f \text{ es continua en } I \iff f \text{ es continua en todo punto de } I.
\]

Desplegada, esta definición lleva \textbf{una cuantificación más} encima de la
del límite: $\forall x_0 \in I,\ \forall \varepsilon > 0,\ \exists \delta > 0,\
\forall x \in I, \dots$ Cuatro cuantificadores anidados.

\begin{notebox}[Para qué tanta maquinaria]
Todo el interés del capítulo es verificar que esta definición corresponde a la
idea superintuitiva de una función \emph{cuya gráfica es de una sola pieza}.
Hace falta la formalización porque sin ella no hay manera de decidir si una
función complicada es o no continua: la intuición del dibujo no alcanza.
\end{notebox}

\section{Catálogo intuitivo de comportamientos}

Antes de los ejemplos formales, el docente hace dibujos, porque en general al
ver la gráfica ya se puede \textbf{predecir} si hay límite o no, y eso orienta
el razonamiento posterior. Los casos que aparecen en el pizarrón:

\begin{center}
\begin{tabularx}{\textwidth}{@{}Xll@{}}
\toprule
\textbf{Comportamiento en $x_0$} & \textbf{¿Tiene límite?} & \textbf{¿Es continua?} \\
\midrule
La gráfica pasa por el límite & sí & sí \\
Salto: hay límite pero $f(x_0)$ está en otro lado & sí & no \\
Oscilación infinitamente rápida & no & no \\
Límites laterales distintos & no & no \\
\bottomrule
\end{tabularx}
\end{center}

El ejemplo típico de oscilación es
\[
  f(x) = \begin{cases}
    \sin\!\left(\dfrac{1}{x}\right) & \text{si } x \neq 0, \\[0.8em]
    42 & \text{si } x = 0,
  \end{cases}
\]
donde el valor en $0$ hay que ponerlo para completar la definición y da
exactamente igual cuál sea (el docente elige $42$ ``porque es la respuesta del
universo''). Esta función \textbf{no tiene límite en $0$}, y el problema no
viene del valor asignado en $0$ sino de que la función oscila demasiado rápido
en torno al origen: no hay ningún $L$ al que se acerque.

El ejemplo típico de límites laterales distintos es la \textbf{función signo}:
por la izquierda tiende a $-1$, por la derecha a $+1$, y como ambos límites
laterales difieren, no hay límite en $0$. Es el contraejemplo que se demuestra
formalmente en la sección \ref{sec:signo}.

% Clase 19 --- Catalogo intuitivo de comportamientos en un punto (§3).
% El docente hace estos cuatro dibujos antes de los ejemplos formales: "en
% general, cuando ven la gráfica, ya pueden predecir si tiene límite o si no
% tiene límite". Incluye sen(1/x) ("una función que oscila demasiado rápido")
% y la función signo ("ambos límites laterales son distintos").
\begin{center}
\begin{tikzpicture}
\begin{axis}[calcfig, width=6.4cm, height=4.2cm, axis lines=middle,
    xmin=-1.6, xmax=1.6, ymin=-0.6, ymax=2.15, xtick=\empty, ytick=\empty,
    xlabel={}, ylabel={}, title={\small (a) continua},
    title style={font=\small, yshift=-2pt}]
  \addplot[domain=-1.5:1.5, samples=60] {0.9 + 0.45*x};
  \node[figptocerrado] at (axis cs:0,0.9) {};
\end{axis}
\end{tikzpicture}\hspace{0.6em}
\begin{tikzpicture}
\begin{axis}[calcfig, width=6.4cm, height=4.2cm, axis lines=middle,
    xmin=-1.6, xmax=1.6, ymin=-0.6, ymax=2.15, xtick=\empty, ytick=\empty,
    xlabel={}, ylabel={}, title={\small (b) salto: $L \neq f(x_0)$},
    title style={font=\small, yshift=-2pt}]
  \addplot[domain=-1.5:-0.03, samples=40] {0.9 + 0.45*x};
  \addplot[domain=0.03:1.5, samples=40, forget plot] {0.9 + 0.45*x};
  \node[figptoabierto] at (axis cs:0,0.9) {};
  \node[figpto, fill=figred, minimum size=5pt] at (axis cs:0,1.62) {};
\end{axis}
\end{tikzpicture}\\[1em]
\begin{tikzpicture}
\begin{axis}[calcfig, width=6.4cm, height=4.2cm, axis lines=middle,
    xmin=-0.62, xmax=0.62, ymin=-1.5, ymax=1.5, xtick=\empty, ytick=\empty,
    xlabel={}, ylabel={}, title={\small (c) oscila: $\sin(1/x)$},
    title style={font=\small, yshift=-2pt}]
  \addplot[domain=0.028:0.6, samples=900] {sin(deg(1/x))};
  \addplot[domain=-0.6:-0.028, samples=900, forget plot] {sin(deg(1/x))};
\end{axis}
\end{tikzpicture}\hspace{0.6em}
\begin{tikzpicture}
\begin{axis}[calcfig, width=6.4cm, height=4.2cm, axis lines=middle,
    xmin=-1.6, xmax=1.6, ymin=-1.6, ymax=1.6, xtick=\empty, ytick=\empty,
    xlabel={}, ylabel={}, title={\small (d) laterales distintos},
    title style={font=\small, yshift=-2pt}]
  \addplot[domain=-1.5:-0.02, samples=2] {-1};
  \addplot[domain=0.02:1.5, samples=2, forget plot] {1};
  \node[figptoabierto] at (axis cs:0,-1) {};
  \node[figptoabierto] at (axis cs:0,1) {};
  \node[figptocerrado] at (axis cs:0,0) {};
\end{axis}
\end{tikzpicture}\\[0.5em]
{\small\itshape Los cuatro casos del pizarrón. En (a) hay límite y coincide
con el valor: es continua. En (b) hay límite pero $f(x_0)$ está en otro lado.
En (c) no hay límite en $0$ porque la función oscila cada vez más rápido (el
valor que se le asigne en $0$ es irrelevante). En (d) los dos límites laterales
existen pero difieren, así que tampoco hay límite.}
\end{center}


\section{Ejemplo 1: la función identidad}

Empieza la parte formal, con la función más sencilla después de las
constantes.

\begin{keybox}[Proposición: la identidad es continua]
Sea $f : \mathbb{R} \to \mathbb{R}$, $f(x) = x$. Entonces
\[
  \forall x_0 \in \mathbb{R}, \quad \lim_{x \to x_0} f(x) = x_0 = f(x_0),
\]
es decir, $f$ es continua en todo punto de $\mathbb{R}$, y por lo tanto
continua en $\mathbb{R}$.
\end{keybox}

\textbf{Demostración.} Sea $\varepsilon > 0$ (lo da el usuario). Se
\textbf{define}
\[
  \boxed{\delta_\varepsilon := \varepsilon}
\]
---por simetría del dibujo: la banda de precisión y la de radio son la
misma---. Sea ahora $x \in \mathbb{R}$ tal que
$0 < |x - x_0| < \delta_\varepsilon$. La desigualdad izquierda no se usa para
nada, así que se descarta. Por definición de $\delta_\varepsilon$,
\[
  |x - x_0| < \varepsilon,
\]
y como $f(x) = x$ por definición de $f$, esto es exactamente
\[
  |f(x) - x_0| < \varepsilon. \qquad \blacksquare
\]

\begin{notebox}[Dónde se usa la definición de $f$]
En el último paso, al reemplazar $x$ por $f(x)$. En toda demostración de
límite hay que usar en algún momento la definición de la función; acá, como la
función es trivial, el razonamiento también lo es. Nótese también que el
límite $L$ vale $x_0$: una vez fijado $x_0$, el valor del límite ya está
determinado y \textbf{no depende de $x$}.
\end{notebox}

\section{Ejercicio: la identidad modificada}

Propuesto en clase para hacer en casa, con las ideas discutidas en el
pizarrón. Sea
\[
  f(x) = \begin{cases}
    x & \text{si } x \neq 0, \\
    42 & \text{si } x = 0.
  \end{cases}
\]

Se plantean tres preguntas: ¿tiene límite en $0$? (sí, vale $0$); ¿tiene
límite en un $x_0 \neq 0$? (sí, vale $x_0$); ¿qué se deduce sobre la
continuidad? La función tiene límite en \textbf{todo} punto y ese límite vale
siempre $x_0$; pero en $x_0 = 0$ el límite $0$ no coincide con el valor
$f(0) = 42$. Conclusión:
\[
  \boxed{f \text{ es continua en todo punto salvo en } 0,
  \text{ donde hay discontinuidad del segundo tipo.}}
\]

\subsection{Cambiar la función en finitos puntos no cambia los límites}

El ejercicio ilustra una propiedad general importante: si se modifica una
función en una \textbf{cantidad finita de puntos}, sus propiedades de límite no
cambian. La demostración se parte en dos casos, y son de dificultad muy
distinta:

\begin{itemize}
  \item \textbf{$x_0 = 0$}: no hay ningún trabajo extra. Como la definición usa
        un entorno \textbf{reducido}, el punto $0$ está excluido y el valor
        $42$ nunca se mira. Sirve el mismo $\delta_\varepsilon = \varepsilon$
        que en la identidad.
  \item \textbf{$x_0 \neq 0$}: acá sí hay que trabajar un poco, porque si
        $\varepsilon$ es grande el entorno de radio $\varepsilon$ alrededor de
        $x_0$ puede \textbf{contener al $0$}, y ahí adentro la función salta a
        $42$. Hay que achicar el radio para esquivar el punto problemático:
        \[
          \boxed{\delta_\varepsilon := \min\left(\varepsilon,\ |x_0|\right)}
        \]
\end{itemize}

% Clase 19 --- La identidad modificada y el radio recortado (§5).
% El docente: "voy a tomar la identidad si x es distinto de 0, pero en 0 voy a
% cambiar el valor... voy a poner 42... cuando se va a considerar un punto acá,
% habrá que reducir un poquito el delta para evitar este valor problemático...
% tomaremos el mínimo entre epsilon y el valor absoluto de x sub 0".
\begin{center}
\begin{tikzpicture}[x=1.25cm, y=1.25cm]
  % banda horizontal de precisión alrededor de L = x_0
  \fill[figamber!18] (-2.35,1.25) rectangle (3.0,2.55);
  % banda vertical de radio delta = min(epsilon, |x_0|): no llega a tocar el 0
  \fill[figblue!14] (1.25,-0.75) rectangle (2.55,3.0);

  % ejes
  \draw[figaxis] (-2.5,0) -- (3.25,0) node[figlbl, below right] {$x$};
  \draw[figaxis] (0,-2.5) -- (0,3.35) node[figlbl, above left] {$y$};

  % la recta y = x, con el punto 0 excluido
  \draw[figline] (-2.35,-2.35) -- (-0.07,-0.07);
  \draw[figline] (0.07,0.07) -- (2.95,2.95);
  \node[figptoabierto] at (0,0) {};
  \node[figlblaux, anchor=north east] at (-0.06,-0.06) {$0$};

  % corte de escala en el eje y, y el valor problemático f(0) = 42
  \draw[figgray, line width=0.7pt] (-0.11,2.86) -- (0.11,3.02);
  \draw[figgray, line width=0.7pt] (-0.11,3.00) -- (0.11,3.16);
  \node[figpto, fill=figred, minimum size=5.5pt] (p42) at (0,3.28) {};
  \node[figlbl, figred, right=4pt] at (p42) {$f(0)=42$};

  % punto de observación x_0 y su imagen
  \draw[figaux] (1.9,0) -- (1.9,1.9) -- (0,1.9);
  \node[figptocerrado] at (1.9,1.9) {};
  \node[figlbl, anchor=south west] at (1.94,0.06) {$x_0$};
  \node[figlbl, anchor=south east] at (-0.06,1.94) {$x_0$};

  % el radio, medido abajo de todo
  \draw[figred, line width=1pt, {Latex[length=2.2mm,width=1.6mm]}-{Latex[length=2.2mm,width=1.6mm]}]
    (1.25,-0.5) -- (2.55,-0.5);
  \draw[figred, line width=0.7pt] (1.25,-0.62) -- (1.25,-0.38);
  \draw[figred, line width=0.7pt] (2.55,-0.62) -- (2.55,-0.38);
  \node[figlbl, figred, anchor=north] at (1.9,-0.62)
    {$\delta_\varepsilon=\min(\varepsilon,|x_0|)$};
\end{tikzpicture}\\[0.5em]
{\small\itshape La identidad con el valor cambiado en un único punto. En
$x_0=0$ no hay ningún trabajo extra: el entorno reducido excluye el $0$ y el
valor $42$ nunca se mira (de ahí el corte de escala en el eje), así que el
límite vale $0$. En un $x_0 \neq 0$ hay que recortar el radio a
$\min(\varepsilon,|x_0|)$ para que la banda vertical no llegue a tocar el punto
problemático. La función tiene límite en todo punto, pero es discontinua
en $0$.}
\end{center}


\subsection{Achicar el radio: qué significa ``suficientemente pequeño''}

Detrás de ese $\min$ hay un argumento completamente general, que ya se podía
observar en el ejemplo de la identidad.

\begin{keybox}[Si un radio funciona, cualquier radio más chico también]
En la identidad se podía haber tomado $\varepsilon/2$, $\varepsilon/42$ o
$\varepsilon/1000$ en lugar de $\varepsilon$ y la demostración andaba igual.
Eso es exactamente lo que quiere decir la expresión \emph{``existe un radio
suficientemente pequeño''}: existe un radio, y la propiedad se conserva al
reemplazarlo por cualquier otro más chico.
\end{keybox}

Es una libertad que se va a usar sistemáticamente, y en la sección siguiente
aparece también del lado de $\varepsilon$.

\section{Ejemplo 2: la raíz cuadrada es continua}

Ahora, en palabras del docente, un ejemplo ``para adultos''. Sea
\[
  f : [0, +\infty) \to \mathbb{R}, \qquad f(x) = \sqrt{x}.
\]
Nótese que el dominio ya \textbf{no} es todo $\mathbb{R}$: es un intervalo no
trivial, los reales positivos o nulos. La gráfica es la media parábola
conocida.

\begin{keybox}[Proposición: la raíz cuadrada es continua]
\[
  \forall x_0 \in [0,+\infty), \quad
  \lim_{x \to x_0}\sqrt{x} = \sqrt{x_0} = f(x_0).
\]
\end{keybox}

La demostración es \textbf{artesanal} y se parte en dos casos de dificultad muy
desigual.

\subsection{Caso fácil: $x_0 = 0$}

Acá $L = \sqrt{0} = 0$, y hay una simplificación gratis: como la función es
positiva, la mitad inferior de la banda de precisión nunca se usa. La idea
geométrica para elegir el radio es que \textbf{para ir del eje $x$ al eje $y$
se aplica la raíz, así que para ir al revés hay que aplicar el cuadrado}. Se
define entonces
\[
  \boxed{\delta_\varepsilon := \varepsilon^{2} > 0}
\]

Sea $x \in [0,+\infty)$ tal que $0 < |x - 0| < \delta_\varepsilon$. Como $x$ es
positivo el valor absoluto se puede sacar, y queda
\[
  0 < x < \varepsilon^{2}.
\]
Como la raíz cuadrada es \textbf{estrictamente monótona} y los tres números son
positivos o nulos, se le puede aplicar la raíz a toda la cadena:
\[
  0 < \sqrt{x} < \sqrt{\varepsilon^{2}} = \varepsilon.
\]
Y si $\sqrt x$ está entre $0$ y $\varepsilon$, entonces
$|\sqrt{x}| < \varepsilon$; escribiendo $|\sqrt{x} - 0| < \varepsilon$
---restar el límite $0$ no cuesta nada--- se tiene exactamente la definición
del límite. $\blacksquare$

\subsection{Caso difícil: $x_0 > 0$}

\textbf{El objetivo, y el razonamiento hacia atrás.} Lo que se quiere obtener
es
\[
  \sqrt{x_0} - \varepsilon < \sqrt{x} < \sqrt{x_0} + \varepsilon,
\]
y de ahí hay que \textbf{extraer} cuál es el $\delta$ que lo garantiza. Lo
natural es elevar al cuadrado para deshacerse de las raíces. Pero elevar al
cuadrado sólo preserva desigualdades \textbf{entre números positivos o nulos},
y ahí aparece el problema: $\sqrt{x_0} + \varepsilon$ es positivo sin
discusión, y $x$ también, pero \textbf{$\sqrt{x_0} - \varepsilon$ puede ser
negativo}, porque nada le prohíbe al usuario pedir una precisión enorme (¿por
qué no medir una raíz cuadrada con precisión de un millón?).

\textbf{La reducción de precisión.} La salida es achicar $\varepsilon$, usando
la libertad de la sección anterior aplicada del otro lado: si el razonamiento
funciona para una precisión más fina, con más razón funciona para la original,
porque la banda más chica está contenida en la más grande. Se define
\[
  \boxed{\varepsilon' := \min\left(\varepsilon,\ \sqrt{x_0}\right) > 0}
\]
que es positivo porque $\varepsilon > 0$ y $x_0 > 0$ por hipótesis. El $\min$
---no el promedio, como se preguntó en clase--- cumple \textbf{dos} cosas a la
vez, y las dos hacen falta: $\varepsilon' \leq \sqrt{x_0}$ garantiza que
$\sqrt{x_0} - \varepsilon' \geq 0$, y $\varepsilon' \leq \varepsilon$ es lo que
va a permitir concluir al final.

\begin{warnbox}[La precisión se achica, nunca se agranda]
Reemplazar $\varepsilon$ por un $\varepsilon' \leq \varepsilon$ es lícito
porque alcanzar la banda más chica implica estar en la más grande. Al revés no
vale. Y la equivalencia con los cuadrados que viene a continuación sólo es
válida porque, después de la reducción, \textbf{todos} los números en juego son
positivos o nulos.
\end{warnbox}

Con $\varepsilon'$ vale entonces la \textbf{equivalencia lógica}
\[
  \sqrt{x_0} - \varepsilon' < \sqrt{x} < \sqrt{x_0} + \varepsilon'
  \iff
  \left(\sqrt{x_0} - \varepsilon'\right)^{2} < x
  < \left(\sqrt{x_0} + \varepsilon'\right)^{2}.
\]

\textbf{El radio.} El objetivo quedó expresado como un intervalo alrededor de
$x$ que \textbf{no es simétrico} respecto de $x_0$: las distancias de $x_0$ a
cada extremo son distintas. Como el radio de seguridad tiene que caber en ambos
lados, se toma el mínimo:
\[
  \boxed{\ \delta_\varepsilon := \min\Big\{
    \left(\sqrt{x_0} + \varepsilon'\right)^{2} - x_0,\ \
    x_0 - \left(\sqrt{x_0} - \varepsilon'\right)^{2} \Big\}\ }
\]

% Clase 19 --- La raíz cuadrada: el radio asimetrico del caso x_0 > 0 (§6.2).
% El docente: "acá voy a tener una distancia entre este número y la raíz de X
% sub cero, y la distancia entre este número y la raíz de X sub cero. Si yo
% tomo el mínimo, esto me va a dar el radio de seguridad." La precisión se
% achica antes a epsilon' = min(epsilon, raíz de x_0) para que
% raíz(x_0) - epsilon' no sea negativo.
\begin{center}
\begin{tikzpicture}
\begin{axis}[
    calcfig,
    axis lines=left,
    width=12cm, height=7cm,
    xmin=0, xmax=3.2, ymin=0, ymax=1.95,
    xtick=\empty, ytick={0.9149,1.2649,1.6149},
    yticklabels={$\sqrt{x_0}-\varepsilon'$, $\sqrt{x_0}$, $\sqrt{x_0}+\varepsilon'$},
    yticklabel style={font=\footnotesize},
    xlabel={$x$}, ylabel={$y$},
    xlabel style={font=\small, at={(axis description cs:1,0)}, anchor=north},
    ylabel style={font=\small, at={(axis description cs:0,1)}, anchor=east},
  ]
  % banda horizontal de precision (ya achicada a epsilon')
  \fill[figamber!18] (axis cs:0,0.9149) rectangle (axis cs:3.05,1.6149);
  % radio de seguridad: el minimo de las dos distancias, centrado en x_0
  \fill[figblue!16] (axis cs:0.8371,0) rectangle (axis cs:2.3629,1.95);

  \addplot[domain=0:3.2, samples=120, figblue, line width=1.1pt] {sqrt(x)};

  % los dos cuadrados que recorta la banda de precision
  \draw[figaux] (axis cs:0.8371,0) -- (axis cs:0.8371,0.9149);
  \draw[figaux] (axis cs:2.6079,0) -- (axis cs:2.6079,1.6149);
  \draw[figaux] (axis cs:1.6,0) -- (axis cs:1.6,1.2649);

  % las dos distancias, en alturas distintas para que no choquen
  \draw[figvec] (axis cs:1.6,0.30) -- (axis cs:0.8371,0.30);
  \node[figlbl, figred, anchor=south] at (axis cs:1.20,0.33) {$d_-$};
  \draw[figvec] (axis cs:1.6,0.62) -- (axis cs:2.6079,0.62);
  \node[figlbl, figred, anchor=south] at (axis cs:2.11,0.65) {$d_+$};

  \node[figlbl, anchor=north] at (axis cs:1.6,-0.16) {$x_0$};
  \node[figlblaux, anchor=north, align=center] at (axis cs:0.8371,-0.03)
    {$(\sqrt{x_0}-\varepsilon')^2$};
  \node[figlblaux, anchor=north, align=center] at (axis cs:2.6079,-0.03)
    {$(\sqrt{x_0}+\varepsilon')^2$};
  \node[figlblaux, anchor=south, align=center] at (axis cs:1.6,1.68)
    {$\delta_\varepsilon = \min(d_-,d_+)$};
\end{axis}
\end{tikzpicture}\\[0.5em]
{\small\itshape La banda de precisión $\sqrt{x_0}\pm\varepsilon'$ recorta
sobre el eje $x$ un intervalo que \emph{no} es simétrico respecto de $x_0$: las
dos distancias $d_-=x_0-(\sqrt{x_0}-\varepsilon')^2$ y
$d_+=(\sqrt{x_0}+\varepsilon')^2-x_0$ son distintas. Como el radio tiene que
caber de los dos lados, se toma $\delta_\varepsilon=\min(d_-,d_+)$ --- de ahí
viene la fórmula horrible de la demostración.}
\end{center}


\textbf{Hay que verificar que este número es positivo}, y no es obvio a la
vista. El argumento es económico: la equivalencia de arriba vale para cualquier
número, en particular para $x = x_0$, y ahí el lado izquierdo es trivialmente
cierto ($\sqrt{x_0}$ está entre $\sqrt{x_0}-\varepsilon'$ y
$\sqrt{x_0}+\varepsilon'$). Luego vale el lado derecho, o sea que $x_0$ está
estrictamente entre los dos cuadrados, y por lo tanto \textbf{ambas diferencias
son positivas} --- y su mínimo también.

\textbf{Verificación.} Sea $x \in [0,+\infty)$ con
$0 < |x - x_0| < \delta_\varepsilon$. Esa desigualdad dice que $x$ está entre
$x_0 - \delta_\varepsilon$ y $x_0 + \delta_\varepsilon$. Se acota cada lado por
separado usando que $\delta_\varepsilon$ es el mínimo:

\begin{itemize}
  \item Por arriba:
    $x < x_0 + \delta_\varepsilon
      \leq x_0 + \left[\left(\sqrt{x_0}+\varepsilon'\right)^{2} - x_0\right]
      = \left(\sqrt{x_0}+\varepsilon'\right)^{2}$.
  \item Por abajo:
    $\delta_\varepsilon \leq x_0 - \left(\sqrt{x_0}-\varepsilon'\right)^{2}$, y
    al cambiar de signo se \textbf{invierte} la desigualdad, de donde
    $x > x_0 - \delta_\varepsilon \geq \left(\sqrt{x_0}-\varepsilon'\right)^{2}$.
\end{itemize}

Pegando ambas:
\[
  \left(\sqrt{x_0} - \varepsilon'\right)^{2} < x
  < \left(\sqrt{x_0} + \varepsilon'\right)^{2}.
\]
Los tres números son positivos o nulos, así que se les puede tomar raíz
cuadrada, y por la equivalencia de antes:
\[
  \sqrt{x_0} - \varepsilon' < \sqrt{x} < \sqrt{x_0} + \varepsilon'
  \quad\Longrightarrow\quad
  \left|\sqrt{x} - \sqrt{x_0}\right| < \varepsilon' \leq \varepsilon,
\]
que es exactamente $|f(x) - L| < \varepsilon$. $\blacksquare$

\subsection{Moraleja: para eso están los teoremas}

\begin{notebox}[De dónde salió ese $\delta_\varepsilon$]
El $\delta_\varepsilon$ horrible de arriba \textbf{fue construido únicamente
para que este razonamiento cerrara}. No es intuitivo, no se adivina: sale de
hacer el razonamiento al revés desde el objetivo.
\end{notebox}

El docente es explícito sobre por qué hizo esto: \emph{``es la última vez que
voy a hacer las cosas a mano''}. La razón de que existan tantos teoremas sobre
límites es que demostrarlos a mano es difícil incluso para funciones sencillas
---con la función cuadrado el trabajo sería parecido---, y hacer esto para cada
función sería inviable. Pero hay que hacerlo \textbf{una vez en la vida}: para
justificar que los teoremas hacen falta, primero hay que sufrir su ausencia. A
partir de la clase siguiente vienen las herramientas que evitan el trabajo
artesanal.

\section{Negar la definición: qué significa no tener límite}

Hasta acá sólo se vieron ejemplos \textbf{positivos}. Para los contraejemplos
hace falta negar la fórmula, aplicando las reglas usuales de negación de
cuantificadores. Partiendo de
\[
  f(x) \to L \ \text{ cuando } x \to x_0
  \iff \forall \varepsilon > 0,\ \exists \delta > 0,\
  \forall x \in E^{*}_{\delta}(x_0),\ f(x) \in E_{\varepsilon}(L),
\]
la negación intercambia cada $\forall$ por un $\exists$ y viceversa:

\begin{keybox}[Negación: $f$ no tiende a $L$]
\[
  \boxed{\ f(x) \not\to L
  \iff \exists \varepsilon > 0,\ \forall \delta > 0,\
  \exists x \in E^{*}_{\delta}(x_0),\quad |f(x) - L| \geq \varepsilon\ }
\]
\end{keybox}

En palabras: \textbf{existe una precisión que fracasa}, en el sentido de que
para cualquier radio que se proponga hay un \textbf{contraejemplo} --- un punto
del entorno reducido cuya imagen se sale de la banda.

Pero eso niega que $L$ sea \emph{el} límite. Para afirmar que la función
\textbf{no tiene límite} hay que negarlo para todos los candidatos a la vez:
\[
  f \text{ no tiene límite en } x_0
  \iff \forall L \in \mathbb{R},\ \exists \varepsilon > 0,\
  \forall \delta > 0,\ \exists x \in E^{*}_{\delta}(x_0),\
  |f(x) - L| \geq \varepsilon.
\]

Son \textbf{cuatro niveles de cuantificación}. Leída como juego: para todo
intento de límite que proponga el usuario, hay que exhibir una precisión que no
funciona para ningún radio.

\begin{notebox}[¿Por qué $\geq$ y no $>$?]
Pregunta de la clase. Porque la negación de ``menor estricto'' es ``mayor o
igual''. De todos modos, si en un caso concreto se consigue la desigualdad
estricta, mejor --- \emph{lo que puede más puede menos}.
\end{notebox}

\section{Contraejemplo: la función signo no tiene límite en 0}
\label{sec:signo}

\[
  \operatorname{sgn} : \mathbb{R} \to \mathbb{R}, \qquad
  \operatorname{sgn}(x) = \begin{cases}
    1 & \text{si } x > 0, \\
    0 & \text{si } x = 0, \\
    -1 & \text{si } x < 0.
  \end{cases}
\]

\begin{keybox}[Proposición]
$\operatorname{sgn}$ no tiene límite en $x_0 = 0$. En todo punto negativo tiene
límite $-1$ y en todo punto positivo tiene límite $1$, de modo que es
\textbf{continua en todo $\mathbb{R}$ salvo en $0$}, donde ni siquiera hay
límite.
\end{keybox}

Lo que hay que demostrar es la fórmula de la sección anterior instanciada en
esta función y en $x_0 = 0$: para todo $L \in \mathbb{R}$ existe
$\varepsilon > 0$ tal que para todo $\delta > 0$ existe $x$ con
\[
  0 < |x| < \delta
  \quad \textbf{y} \quad
  \left|\operatorname{sgn}(x) - L\right| \geq \varepsilon.
\]

\begin{notebox}[Ese ``y'' se lee ``pero'']
En matemática el símbolo $\wedge$ se lee indistintamente ``y'' o ``pero'';
sólo cambia la intención, no el significado.
\end{notebox}

\subsection{La estrategia: $\varepsilon$ es la mitad del salto}

El orden del trabajo lo fija la alternancia de cuantificadores: se recibe un
$L$ cualquiera, se \textbf{construye} un $\varepsilon$, se recibe un $\delta$
cualquiera, se \textbf{construye} un $x$ en función de todo lo anterior, y se
verifica. Igual que antes, demostrar cosas sobre límites es construir
funciones; sólo que ahora se construye un contraejemplo.

\begin{keybox}[El truco, y es general]
Cuando hay dos límites laterales distintos, el $\varepsilon$ que funciona es
\textbf{la mitad de la distancia entre ellos}. Acá los laterales son $-1$ y
$1$, a distancia $2$, así que se toma
\[
  \boxed{\varepsilon := 1}
\]
Ni siquiera es una precisión fina: se va a mostrar que ya con precisión $1$
ningún $L$ sirve.
\end{keybox}

Sea entonces $\delta > 0$ cualquiera. Para construir el $x$ se distinguen dos
casos según el signo de $L$; la idea en ambos es \textbf{irse hacia el lado
opuesto a $L$}, donde el valor de la función queda lejos.

% Clase 19 --- El contraejemplo de la función signo (§8).
% El docente: "como hay una distancia de 2 hay que tomar la mitad... se toma
% epsilon igual a 1... si yo quiero mostrar que este número L no es límite,
% para obtener un contraejemplo yo tengo que ir muy lejos de L. Voy a tomar un
% punto negativo, porque el valor asociado será menos uno, que estará lejos
% de L."
\begin{center}
\begin{tikzpicture}[x=2.0cm, y=1.25cm]
  % banda de precision L +- epsilon, con epsilon = 1
  \fill[figamber!18] (-1.75,-0.4) rectangle (1.75,1.6);
  % banda del radio delta alrededor de 0
  \fill[figblue!14] (-0.7,-1.75) rectangle (0.7,1.9);

  % ejes
  \draw[figaxis] (-1.85,0) -- (1.95,0) node[figlbl, below right] {$x$};
  \draw[figaxis] (0,-1.85) -- (0,2.1) node[figlbl, above left] {$y$};

  % la funcion signo
  \draw[figline] (-1.75,-1) -- (-0.04,-1);
  \draw[figline] (0.04,1) -- (1.75,1);
  \node[figptoabierto] at (0,-1) {};
  \node[figptoabierto] at (0,1) {};
  \node[figptocerrado] at (0,0) {};
  \node[figlbl, right=3pt] at (1.75,1) {$1$};
  \node[figlbl, left=3pt] at (-1.75,-1) {$-1$};

  % el intento de limite L (positivo o nulo) y su banda
  \draw[figred, dashed, line width=0.8pt] (-1.85,0.6) -- (1.95,0.6);
  \node[figlbl, figred, right=2pt] at (1.95,0.6) {$L$};
  \node[figlblaux, figamber, anchor=west] at (1.78,1.6) {$L+\varepsilon$};
  \node[figlblaux, figamber, anchor=west] at (1.78,-0.4) {$L-\varepsilon$};

  % el contraejemplo x = -delta/2
  \node[figpto, fill=figred, minimum size=5.5pt] (cx) at (-0.35,-1) {};
  \draw[figaux] (-0.35,0) -- (cx);
  \node[figlbl, figred, anchor=north, align=center] at (-0.35,-1.15)
    {$x=-\dfrac{\delta}{2}$};
  \node[figlblaux, anchor=south] at (-0.7,1.92) {$-\delta$};
  \node[figlblaux, anchor=south] at (0.7,1.92) {$\delta$};
\end{tikzpicture}\\[0.5em]
{\small\itshape Sea cual sea el intento de límite $L$, la precisión
$\varepsilon=1$ --- la mitad de la distancia entre los dos límites laterales ---
fracasa. Para $L\geq 0$ basta ir al lado negativo, donde el signo vale $-1$: el
punto $x=-\delta/2$ está en el entorno reducido de radio $\delta$ pero su
imagen queda fuera de la banda $L\pm 1$, ya que $|-1-L|=L+1\geq 1$. Para $L<0$
se toma $x=\delta/2$ y el argumento es simétrico.}
\end{center}


\subsection{Caso $L \geq 0$}

Si $L$ es positivo o nulo, hay que ir a la parte \textbf{negativa}, donde el
signo vale $-1$. Se toma
\[
  x := -\frac{\delta}{2},
\]
que cumple la condición del entorno reducido por construcción, ya que
$0 < |x| = \delta/2 < \delta$. Como $x < 0$, se tiene
$\operatorname{sgn}(x) = -1$, y entonces
\[
  \left|\operatorname{sgn}(x) - L\right| = \left|-1 - L\right|
  = \left|L + 1\right| = L + 1 \geq 1 = \varepsilon,
\]
donde el valor absoluto se saca porque $L + 1$ es positivo (al ser $L \geq 0$).

\subsection{Caso $L < 0$}

Simétricamente, se va a la parte \textbf{positiva}. Se toma
\[
  x := \frac{\delta}{2},
\]
que de nuevo cumple $0 < |x| < \delta$. Como $x > 0$, es
$\operatorname{sgn}(x) = 1$, y
\[
  \left|\operatorname{sgn}(x) - L\right| = \left|1 - L\right|
  = 1 - L > 1 = \varepsilon,
\]
porque $L < 0$ implica $-L > 0$, así que $1 - L$ es positivo (se puede sacar el
valor absoluto) y mayor que $1$: \emph{uno menos un número negativo es mayor
que uno}.

En los dos casos se encontró el contraejemplo pedido, para cualquier $L$ y
cualquier $\delta$. Luego $\operatorname{sgn}$ no tiene límite en $0$.
$\blacksquare$

\vspace{1em}
\noindent\emph{La próxima clase abandona el trabajo artesanal: se demuestran
las propiedades algebraicas de los límites (suma, producto, cociente), que
permiten calcular límites de funciones complicadas sin volver nunca a la
definición $\varepsilon$-$\delta$.}

\end{document}
